\documentclass[ocean-paper.tex]{subfiles}
\begin{document}
\section{Observation Space}\label{appendix:observationspace}
\begin{figure}
    \centering
    \includegraphics[width=0.75\textwidth]{figures/dotahud.pdf}
    \caption{\dota's human ``Observation Space''}
    \label{fig:dotahud}
\end{figure}
\begin{figure}
    \centering
    \includegraphics[width=\textwidth]{figures/unit-observations.png}
    \caption{\textbf{Observation Space Overview:} 
    The arrays that \openaifive observes at each timestep.
    Most of \openaifive's observations  are unit-centered; for 189 different units on the map, we observe a set of basic properties. These units are grouped along the top of the figure.
    We observe some data about all units, some extra data about the primary units (the heroes), and even more data about the heroes on our team.
    A few observations are not tied to any unit. 
    Finally, two observations having to do with hero control (terrain near me, and my previous action) are only observed about the individual hero that this LSTM replica operates.
    In this diagram blue bands represent categorical data and yellow bands represent continuous or boolean data; most entities (units, modifiers, abilities, items, and pickups), have some of each.
    Each piece of the figure summarizes the total dimensionality of that portion of the input.
    All together, an \openaifive hero observes 1,200 categorical values % 10 + 10 + 164*4 + 30*2 + 10 + 100 + 60 + 160 + 6 + 128
    and 14,534 continuous/boolean % 215*2 + 3526*2 + 645 + 125*2 + 1055 + 100*2 + 328*2 + 60 + 210*2 + 1040*2 + 22 + 90 + 900 + 384 + 310
    values.\nicetohave{Add little images to the minimap and other places where it will help.}}
    \label{fig:obs-space}
\end{figure}


\begin{table}
%%% This table's formatting is really hacky. If anyone wants to fix it that would be great.
\begin{scriptsize}
\begin{threeparttable}
\begin{subtable}[t]{0.3\textwidth}
\begin{tabular}[t]{|p{4cm}|c|}
  \hline
  \textbf{Global data}   & \textbf{22}\\\hline
  time since game started & 1 \\\hline
  is it day or night? &  \\
  time to next day/night change & 2 \\\hline
  time to next spawn: creep, neutral, bounty, runes & 4 \\\hline
  time since seen enemy courier &  \\
  is that $>$ 40 seconds?\tnote{a} & 2 \\\hline
  min\&max time to Rosh spawn & 2 \\\hline
  Roshan's current max hp & 1 \\\hline
  is Roshan definitely alive? & 1 \\\hline
  is Roshan definitely dead? & 1 \\\hline
  Next Roshan drops cheese? & 1 \\\hline
  Next Roshan drops refresher? & 1 \\\hline
  Roshan health randomization\tnote{b} & 1 \\\hline
  Glyph cooldown (both teams) & 2 \\\hline
  Stock counts\tnote{c} & 4 \\\hline
\end{tabular}
\begin{tabular}{|p{4cm}|c|}
  \hline
  \textbf{Per-unit (189 units)}   &  \textbf{43}\\\hline
  position (x, y, z) & 3 \\\hline
  facing angle (cos, sin) & 2 \\\hline
  currently attacking?\tnote{e} & \\
  time since last attack\tnote{d} & 2 \\\hline
  max health & \\
  last 16 timesteps' hit points & 17 \\\hline
  attack damage, attack speed & 2 \\\hline
  physical resistance & 1 \\\hline
  invulnerable due to glyph? & \\
  glyph timer & 2 \\\hline
  movement speed & 1 \\\hline
  on my team? neutral? & 2 \\\hline
  animation cycle time & 1 \\\hline
  eta of incoming ranged \& tower creep projectile (if any) & \\
  \# melee creeps atking this unit\tnote{d} & 3 \\\hline
  [Shrine only] shrine cooldown & 1 \\\hline
  vector to me (dx, dy, length)\tnote{e} & 3 \\\hline
  am I attacking this unit?\tnote{e} & \\
  is this unit attacking me?\tnote{d,e} & \\
  eta projectile from unit to me\tnote{e} & 3 \\\hline
  % \hline
  % textbf{Per-unit categorical data}   &  \\\hline
  \categorical unit type & \categorical 1 \\\hline
  \categorical current animation & \categorical 1 \\
  \hline
\end{tabular}
\end{subtable}
\hfill
\begin{subtable}[t]{0.3\textwidth}
\begin{tabular}[t]{|p{4cm}|c|}
  \hline
  \textbf{Per-hero add'l (10 heroes)}   &  \textbf{25}\\\hline
  is currently alive? & 1 \\\hline
  number of deaths & 1 \\\hline
  hero currently in sight? & \\
  time since this hero last seen & 2 \\\hline
  hero currently teleporting? & \\
  if so, target coordinates (x, y) & \\
  time they've been channeling & 4 \\\hline
  respawn time & 1 \\\hline
  current gold (allies only) & 1 \\\hline
  level & 1 \\\hline
  mana: max, current, \& regen & 3 \\\hline
  health regen rate & 1 \\\hline
  magic resistance & 1 \\\hline
  strength, agility, intelligence & 3 \\\hline
  currently invisible? & 1 \\\hline
  % Due to a bug this is always zero.
  % is\_in\_ability\_phase & \\
  % I think this is the same as the next one.
  % has\_active\_ability & \\
  is using ability? & 1\\\hline
  \# allied/enemy creeps/heroes in line btwn me and this hero\tnote{e} & 4 \\\hline
  \hline
  \textbf{Per-allied-hero additional (5 allied heroes)}   & \textbf{211} \\\hline
  Scripted purchasing settings\tnote{b} & 7 \\\hline
  Buyback: has?, cost, cooldown & 3 \\\hline
  Empty inventory \& backpack slots & 2 \\\hline
  Lane Assignments\tnote{b} & 3 \\\hline
  Flattened nearby terrain: 14x14 grid of passable/impassable? & 196 \\\hline
  \categorical scripted build id & \categorical\\
  \categorical next item to purchase\tnote{b} & \categorical 2 \\\hline
  \hline
  \textbf{Nearby map (8x8)\tnote{e}}   & \textbf{6} \\\hline
  terrain: elevation, passable? & 2 \\\hline
  allied \& enemy creep density & 2 \\\hline
  area of effect spells in effect.\tnote{f} & 2 \\\hline
  \categorical area of effect spells in effect.\tnote{f} & \categorical 2 \\\hline
  \hline
  \textbf{Previous Sampled Action\tnote{e}}   & \textbf{310} \\\hline
  Offset? (Regular, Caster, Ward) & 3x2x9 \\\hline
  % Omitting tp head because that can't possibly matter and it's hard to explain.
  Unit Target's Embedding & 128 \\\hline
  Primary Action's Embedding & 128 \\\hline
\end{tabular}
\end{subtable}
\hfill
\begin{subtable}[t]{0.3\textwidth}
\begin{tabular}[t]{|p{4cm}|c|}
  \hline
  \textbf{Per-modifier (10 heroes x 10 modifiers \& 179 non-heroes x 2 modifiers)}   &  \textbf{2}\\\hline
  remaining duration & 1 \\\hline
  stack count & 1 \\\hline
  \categorical modifier name & \categorical 1 \\
  \hline
  \hline
  \textbf{Per-item (10 heroes x 16 items)}   & \textbf{13}\\\hline
  location one-hot (inventory/backpack/stash) & 3 \\\hline
  charges & 1 \\\hline
  is on cooldown? & \\
  cooldown time & 2 \\\hline
  is disabled by recent swap? & \\
  item swap cooldown & 2 \\\hline
  toggled state & 1 \\\hline
  special Power Treads one-hot & \\
  (str/agi/int/none) & 4\\\hline
  \categorical item name & \categorical 1 \\
  \hline
  \hline
  \textbf{Per-ability (10 heroes x 6 abilities)}   &  \textbf{7}\\\hline
  cooldown time & 1 \\\hline
  in use? & 1 \\\hline
  castable & 1 \\\hline
  Level 1/2/3/4 unlocked?\tnote{d} & 4\\\hline
  \categorical ability name & \categorical 1 \\
  \hline
  \hline
  \textbf{Per-pickup (6 pickups)}   &  15\\\hline
  status one-hot (present/not present/unknown) & 3 \\\hline
  location (x, y) & 2 \\\hline
  distance from all 10 heroes & 10 \\\hline
  \categorical pickup name & \categorical 1 \\
  \hline
  \hline
  \textbf{Minimap (10 tiles x 10 tiles)}   &  \textbf{9}\\\hline
  fraction of tile visible & 1 \\\hline
  \# allied \& enemy creeps & 2 \\\hline
  \# allied \& enemy wards & 2 \\\hline
  \# enemy heroes & 1 \\\hline
  cell (x, y, id) & 3 \\\hline
  % "id" is my description of a really confusing thing that appears to be number of neutral camps in the cell (Roshan counts as 2) + 3 if there's a bounty rune.
  % It's actually a 10x10x12 array, but we appear to never use indices 1 and 2, and index 3 is identical to index 4.
\end{tabular}
\end{subtable}
\end{threeparttable}
\begin{threeparttable}
\begin{tablenotes}
    \item[a] These observations are leftover from an early version of Five which played a restricted 1v1 version of the game. They are likely obsolete and not needed, but this was not tested.
    \item[b] These observations are about our per-game randomizations. See \autoref{sec:methods:exploration}.
    \item[c] For items: gem, smoke of deciept, observer ward, infused raindrop.
    \item[d] Observations are not visible per-se, but can be estimated. We use scripted logic to estimate them from visible observations.
    \item[e] These observations (only) are different for the five different heroes on the team.
    \item[f] This observation appears twice, and serves as an example of the difficulties of surgery. Although this is a categorical input, we began by treating it as a float input to save on engineering work (this observation is unlikely to be very important). Later the time came to upgrade it to a properly embedded categorical input, but our surgery tools do not support removing existing observations. Hence we added the new observation, but were forced to leave the deprecated observation as well.
\end{tablenotes}
\end{threeparttable}
\end{scriptsize}
\caption{
\textbf{Full Observation Space: }
All observations \openaifive receives at each time step. Blue rows are categorical data. Entries with a question mark are boolean observations (only take values 0 or 1 but treated as floats otherwise). The bulk of the observations are per-unit observations, observed for each of 189 units: heroes (5), creeps (30), buildings (21), wards (30), and courier (1) for each team, plus 15 neutrals. If the number of visible units in a category is less than the allotted number, the rest are padded with zeroes. If more, we observe only the units closest to allied heroes. Units in fog of war are not observed. When enemy heroes are in fog of war, we reuse the observation from the last time step when the unit was visible.
}
\label{table:full-obs}
\end{table}

At each time step one of our heroes observes \obssize inputs about the game state (mostly real numbers with some integer categorical data as well). 
See \autoref{fig:obs-space} for a schematic outline of our observation space and \autoref{table:full-obs} for a full listing of the observations. 

Instead of using the pixels on the screen, we approximate the information available to a human player in a set of data arrays. 
This approximation is imperfect; there are small pieces of information which humans can gain access to which we have not encoded in the observations. 
On the flip side, while we were careful to ensure that all the information available to the model is also available to a human, the model does get to see \emph{all} the information available simultaneously every time step, whereas a human needs to click into various menus and options to get that data. 
Although these discrepancies are a limitation, we do not believe they meaningfully detract from our ability to benchmark against human players.

Humans observe the game via a rendered screen, depicted in \autoref{fig:dotahud}. 
\openaifive uses a more semantic observation space than this for two reasons: First, because our goal is to study strategic planning and gameplay rather than focus on visual processing. 
Second, it is infeasible for us to render each frame to pixels in all training games; this would multiply the computation resources required for the project manyfold.

All float observations (including booleans which are treated as floats that happen to take values 0 or 1) are normalized before feeding into the neural network. For each observation, we keep a running mean and standard deviation of all data ever observed; at each timestep we subtract the mean and divide by the st dev, clipping the final result to be within (-5, 5). 

\end{document}